% TOP_PROBLEM: {"id": 3162, "title": "The circular embedding conjecture", "category_id": 3, "category": {"id": 3, "name": "graph_theory", "display_name": "Graph Theory", "description": "Problems involving graphs, networks, and their properties.", "slug": "graph-theory", "order_index": 3, "created_at": "2026-07-31T15:26:25.671Z"}, "status": "open", "problem_number": "OPG-138", "set_id": 12}
% FIRST_POSTED: 2026-10-01
% ENTRY_KIND: statement_only
% UnsolvedMath ID 3162; source catalogue code OPG-138
% !TeX program = lualatex
% Definition and sources only; this file is not an LLM solution attempt.
\documentclass[11pt,a4paper]{article}
\usepackage[margin=25mm]{geometry}
\usepackage{fontspec}
\setmainfont{lmroman10-regular.otf}[BoldFont=lmroman10-bold.otf,
 ItalicFont=lmroman10-italic.otf,BoldItalicFont=lmroman10-bolditalic.otf]
\usepackage{amsmath,amssymb,mathtools}
\usepackage{unicode-math}
\setmathfont{latinmodern-math.otf}
\AtBeginDocument{\let\setminus\smallsetminus}
\usepackage{xurl,hyperref}
\hypersetup{colorlinks=true,urlcolor=blue}
\setcounter{secnumdepth}{0}
\setlength{\parindent}{0pt}
\setlength{\parskip}{5pt}
\setlength{\emergencystretch}{3em}
\begin{document}
\section*{The circular embedding conjecture}\label{3162}
{\small UnsolvedMath ID 3162 · OPG-138 · Graph Theory · OpenGarden}
\subsection{Definitions and mathematical statement}
Let $G$ be a finite simple $2$-connected graph, meaning that
it has at least three vertices and remains connected after
deleting any one vertex. A surface here means a compact
connected topological $2$-manifold without boundary. It may be
orientable or nonorientable, and its genus is not fixed in
advance.

An embedding of $G$ in a surface places vertices at distinct
points and realizes edges as arcs whose interiors are mutually
disjoint and avoid the vertices. It is cellular if each
component of the surface minus the embedded graph is an open
disc; these components are its faces. A facial boundary is
usually described by a closed walk, which may in general
repeat vertices or edges.

The circular embedding conjecture asks whether every such $G$
has a cellular embedding in some surface in which the boundary
walk of every face is a simple cycle. Thus each facial boundary
visits each of its vertices once before returning to its start;
the face closes to a disc bounded by an embedded circle.

The surface and rotation of incident edges may be chosen to
suit $G$. No planar embedding or orientable embedding is
required. Each edge has two sides in a surface, so facial
cycles in such an embedding cover every edge twice when faces
are counted with multiplicity. However, the target includes
the surface condition at vertices: gluing discs along an
arbitrary cycle cover need not give manifold neighbourhoods.
\subsection{Short English statement}
Can every 2-connected graph be embedded on some closed surface so each face is a disc bounded by a simple cycle?
\subsection{Sources}
\begin{itemize}
\item[S1] ULAM.AI and the UnsolvedMath Contributors, supplied problems.json, record 3162 (OPG-138), OpenGarden collection. Catalogue identity and source transcription; CC BY 4.0.
\url{https://huggingface.co/datasets/ulamai/UnsolvedMath}.
\item[S2] Open Problem Garden, The circular embedding conjecture. Original problem and discussion; the mathematical formulation below is expanded from the supplied source record.
\url{http://www.openproblemgarden.org/op/the_circular_embedding_conjecture}.
\end{itemize}
\end{document}
